\chapter{Centered Bloch Rapidity Form of the Collatz Branches}
\label{chap:bloch-rapidity-collatz}

\status{PROVED algebraic conjugacy and exact finite-orbit identity. No universal Collatz-convergence, xi zero-location, or RH claim is made.}

The preceding Uroboros construction becomes particularly simple after one further centering of the Riemann coordinate.  Define
\[
u=2x,\qquad s=\frac{u}{1+u},\qquad t=2s-1=\frac{u-1}{u+1}.
\]
If \(u=e^\lambda\), then
\[
\boxed{t=\tanh(\lambda/2).}
\]
Thus \(t\) is a bounded centered coordinate for the logarithmic radial variable.

\section{The half-layer is the centered zero}

At the distinguished half-layer,
\[
 x=\frac12\quad\Longleftrightarrow\quad u=1
 \quad\Longleftrightarrow\quad s=\frac12
 \quad\Longleftrightarrow\quad \boxed{t=0}.
\]
This is an exact coordinate identity.  It does not identify the arithmetic numbers \(1/2\) and \(0\); rather, the half-layer in \(x\) is mapped to the origin of the centered \(t\)-chart.

\section{General dilation as hyperbolic translation}

Under a positive dilation
\[
u\mapsto au,
\]
direct substitution gives
\[
T_a(t)=
\frac{(a-1)+(a+1)t}{(a+1)+(a-1)t}.
\]
Writing
\[
p_a=\frac{a-1}{a+1},
\]
we obtain
\[
\boxed{T_a(t)=\frac{t+p_a}{1+p_at}.}
\]
Because
\[
p_a=\tanh\!\left(\frac{\log a}{2}\right),
\]
composition of dilations corresponds to the hyperbolic addition law
\[
\boxed{p\oplus q=\frac{p+q}{1+pq}.}
\]

\section{The halving branch}

For one halving step, \(a=1/2\), so
\[
p_{1/2}=-\frac13.
\]
Consequently
\[
\boxed{
B(t)=\frac{t-1/3}{1-t/3}
     =\frac{3t-1}{3-t}.
}
\]
Thus the halving branch is an exact hyperbolic/M\"obius translation of logarithmic rapidity \(-\log2\).

\section{The odd Collatz branch}

The odd branch \(x\mapsto3x+1\) was previously written in the \(s\)-coordinate as
\[
O_s(s)=\frac{s+2}{3}.
\]
Since \(t=2s-1\), it becomes simply
\[
\boxed{O_t(t)=\frac{t+2}{3}.}
\]
The complete branchwise conjugacy is therefore
\[
\boxed{
C_t(t)=
\begin{cases}
\dfrac{3t-1}{3-t}, & x\text{ even},\\[0.75em]
\dfrac{t+2}{3}, & x\text{ odd}.
\end{cases}}
\]
No convergence assumption is involved.

\section{The five-halving orbit}

The exact finite chain
\[
16\to8\to4\to2\to1\to\frac12
\]
becomes
\[
\boxed{
\frac{31}{33}
\to\frac{15}{17}
\to\frac79
\to\frac35
\to\frac13
\to0.
}
\]
The starting value is
\[
\frac{31}{33}
=\frac{32-1}{32+1}
=\tanh\!\left(\frac{5\log2}{2}\right).
\]
Five copies of the single-step parameter \(-1/3\) combine exactly to
\[
\boxed{-\frac{31}{33}}.
\]
Equivalently, the scale \(u\mapsto u/32\) acts as
\[
T_{1/32}(t)
=\frac{t-31/33}{1-(31/33)t},
\]
and therefore
\[
\boxed{T_{1/32}(31/33)=0.}
\]
Hence:

\begin{theorem}[SOH-G009]
In the centered coordinate
\[
t=\frac{u-1}{u+1}=2s-1,
\]
the finite chain \(16\to8\to4\to2\to1\to1/2\) is exactly a five-step hyperbolic translation from \(31/33\) to the centered half-layer \(t=0\):
\[
\boxed{B^5(31/33)=0.}
\]
\end{theorem}

\section{Connection to the scale-32 quotient}

The same rational number \(31/33\) occurs in three exact roles:
\begin{enumerate}
\item it is the centered coordinate of \(u=32\), equivalently \(x=16\);
\item it is the M\"obius parameter \(p_{32}\) of the scale-32 translation;
\item its negative is the composed parameter of five halving steps.
\end{enumerate}
Thus the scale appearing in the Uroboros quotient and the selected five-halving orbit are the same hyperbolic translation written in opposite directions.

\section{Proof firewall}

The theorem concerns an exact finite orbit and exact branch conjugacies.  It does not assert that every Collatz trajectory enters this orbit.  It also does not produce a new functional equation for \(\xi\), a zero-location theorem, SOH-G003, or the Riemann Hypothesis.
